\documentclass[a4paper]{article}
\usepackage{amsfonts}
\usepackage{amsmath}
\usepackage{amssymb}
\usepackage{graphicx}     

\begin{document}
  
\noindent \large Auteur: Odia Kibor \\
\noindent \large Studierichting: Informatica\\
\noindent \large Oefeningenreeks 1E nr. 13\\

\medskip

\normalsize


We delen door het product $(z-1)(z+1)(z-2)$, dat graad 3 heeft.
Dus de rest is maximaal van graad 2.\\

$\Rightarrow R(z) = $ 2de graad $\Rightarrow  az^2 + bz + c$\\

$\Rightarrow z^6 - 1 = (z-1)(z+1)(z-2) \cdot Q(z) + az^2 + bz + c$\\

Vul $z = 1$, $-1$ en $2$ in.\\

$\Rightarrow 1^6 - 1 = 0$\\

$\Rightarrow (-1)^6 - 1 = 0$\\

$\Rightarrow 2^6 - 1 = 63$\\

Stelsel:\\


$\left\{
	\begin{array}{l} 
		a + b + c = 0\\
		a - b + c = 0\\
        4a + 2b + c = 63
	\end{array}
\right .
\Rightarrow
\left\{
	\begin{array}{l} 
		2b = 0\\
		a - b + c = 0\\
        4a + 2b + c = 63
	\end{array}
\right .
\Rightarrow
\left\{
	\begin{array}{l} 
		b = 0\\
		a + c = 0\\
        4a + 2b + c = 63
	\end{array}
\right .
\Rightarrow
\left\{
	\begin{array}{l} 
		b = 0\\
		a + c = 0\\
        3a = 63
	\end{array}
\right .
\Rightarrow
\left\{
	\begin{array}{l} 
		b = 0\\
		a + c = 0\\
        a = 21
	\end{array}
\right .
\Rightarrow
\left\{
	\begin{array}{l} 
		a = 21\\
		b = 0\\
        c = -21
	\end{array}
\right .$\\\\

$R(z) = 21z^2 - 21$









\begin{thebibliography}{999}
%\bibitem{a} Referentie
\end{thebibliography}
\end{document} 


